
\subsubsection{6.1 Summary of Findings}

Three sub-hypotheses were experimentally evaluated out of five planned:

\begin{itemize}
\item \textbf{h-e1 (existence, MUST\_WORK gate): Passed.} Flat-MLP exhibited significant permutation sensitivity (delta-rho = 0.160, p = 0.000); NFT showed no significant sensitivity (delta-rho = 4.09 x $10^-6$, p = 0.477).
\end{itemize}

\begin{itemize}
\item \textbf{h-m1 (mechanism, MUST\_WORK gate): Passed.} NFT-base delta-rho = 4.71 x $10^-7$; mediation delta-R-squared = 0.228 >= 0.10 threshold. Six-encoder ablation across 18 runs confirmed encoder ranking.
\end{itemize}

\begin{itemize}
\item \textbf{h-m2 (alternatives, SHOULD\_WORK gate): Failed.} Augmentation provided partial compensation (67\% delta-rho reduction) but with high seed variance. L2 canonicalization collapsed to constant predictor. The three-way ranking condition could not be evaluated.
\end{itemize}

\begin{itemize}
\item \textbf{h-m3 (graceful degradation) and h-m4 (cross-pipeline transfer): Not executed.}
\end{itemize}

The results support the claim that NFT encoders exhibit substantially lower permutation sensitivity than flat-MLP encoders on this benchmark, and that the difference is mediated by equivariant attention (as measured by the delta-R-squared analysis). The results do not address cross-pipeline transfer robustness, which was part of the original hypothesis but was not experimentally tested.

\subsubsection{6.2 Interpretation of Mediation Analysis}

The delta-R-squared = 0.228 indicates that NFT-base embeddings explained 22.8 percentage points more variance in generalization gap than flat-MLP+aug embeddings. This is interpreted as evidence that equivariant attention captures structural signals (attributed to neuron influence concentration) that augmentation-based approaches do not encode invariantly. However, this analysis compares two architectures of different sizes (75K vs. 3.04M parameters), and the mediation effect could partly reflect architectural differences unrelated to equivariance. A more controlled mediation analysis would require matched-parameter encoders.

\subsubsection{6.3 Flat-MLP Degradation Across Experiments}

Flat-MLP delta-rho was 0.160 in h-e1 (50 epochs) and 0.640 in h-m1 (100 epochs). The larger degradation under longer training is consistent with the hypothesis that extended training increases reliance on neuron ordering statistics, though other differences between experiments (number of seeds, minor implementation variations) could also contribute. Both values exceeded the pre-specified 0.10 threshold.

\subsubsection{6.4 Limitations}

\textbf{Dataset adaptation.} The intended Unterthiner FC-MLP zoo was unavailable (URL returned HTTP 404). The CNN zoo (29,997 models, 4-layer CNNs) was adapted by reshaping weight matrices to per-neuron token format. While the permutation structure is preserved by this adaptation, absolute metric values may differ on native FC-MLP weights, and the generalization gap characteristics may differ between CNN and FC-MLP model families.

\textbf{Cross-pipeline transfer not tested.} The original hypothesis included a prediction about cross-pipeline transfer robustness (MNIST to CIFAR). Sub-hypotheses h-m3 and h-m4 were not executed. No claims about cross-distribution generalization can be made from this study.

\textbf{Canonicalization scope.} Only L2-norm canonicalization was tested as a post-hoc approach. Stronger canonicalization methods (sort-by-magnitude, Hungarian alignment without oracle access, spectral normalization) were not evaluated. The oracle canonicalization result (delta-rho = 0.000) demonstrates that perfect post-hoc alignment is theoretically possible; whether practical non-oracle canonicalization can approach this bound is an open question.

\textbf{Augmentation seed count.} The augmentation analysis is based on 3 seeds. The observed high variance (delta-rho range: 0.096--0.317) is suggestive of a multi-modal optimization landscape but cannot be characterized precisely with so few samples.

\textbf{Unmatched parameter comparison.} NFT (75K parameters) and flat-MLP (3.04M parameters) differ by 40x in parameter count. The observed baseline performance difference (rho = 0.489 vs. 0.303) cannot be attributed solely to equivariance without a matched-parameter control. Flat-MLP may be over-parameterized for a 30K-model zoo (approximately 127 parameters per training example vs. 3 for NFT).

\textbf{Single zoo, single task.} All experiments use a single model zoo (Unterthiner MNIST) and a single prediction target (generalization gap). Generalization to other zoos, architectures, or prediction targets is not established.

\subsubsection{6.5 Broader Considerations}

The mediation analysis framework (delta-R-squared as a test of architectural inductive bias) may be applicable to other settings where encoder architecture is compared against the symmetry structure of the input domain. The finding that L2-norm canonicalization collapses the predictor provides a concrete negative result for weight-space analysis: magnitude-destructive normalization removes information relevant to scalar property prediction from weight vectors.
